% Generated draft: requires mathematical and source review.
\documentclass[11pt]{article}
\usepackage[margin=25mm]{geometry}
\usepackage{fontspec}
\setmainfont{Latin Modern Roman}
\usepackage{amsmath,amssymb,mathtools,unicode-math}
\setmathfont{Latin Modern Math}
\usepackage{enumitem,needspace,xurl,hyperref}
\setlength{\parindent}{0pt}
\setlength{\parskip}{5pt}
\setlength{\emergencystretch}{3em}
\newcommand{\N}{\mathbb{N}}
\newcommand{\Z}{\mathbb{Z}}
\newcommand{\Q}{\mathbb{Q}}
\newcommand{\R}{\mathbb{R}}
\newcommand{\C}{\mathbb{C}}
\newcommand{\F}{\mathbb{F}}
\newcommand{\A}{\mathbb{A}}
\newcommand{\PP}{\mathbb{P}}
\newcommand{\E}{\mathbb{E}}
\newcommand{\eps}{\varepsilon}
\newcommand{\dd}{\,\mathrm{d}}
\begin{document}
\section*{UnsolvedMath ID: 30002691}
\textbf{Source problem code:} catalogue entry
\subsection{Definitions and mathematical statement}
The supplied record contains two questions, with unresolved choices described below.

For the first question, let \(\mathfrak{s}\) be a finite-dimensional real vector space with an inner product, and let \(\mu_0\in\Lambda^2\mathfrak{s}^*\otimes\mathfrak{s}\) be a solvable Lie bracket representing a left-invariant Riemannian metric. Write \(\operatorname{Ric}_{\mu}\) for the Ricci endomorphism determined by the bracket and inner product. For an endomorphism \(A\), put
\[
(\pi(A)\mu)(u,v)=A\mu(u,v)-\mu(Au,v)-\mu(u,Av).
\]
The unnormalized bracket flow is
\[
\frac{\dd\mu}{\dd t}=-\pi(\operatorname{Ric}_{\mu})\mu,
\qquad \mu(0)=\mu_0.
\]
A Riemannian metric \(g\) is a Ricci soliton if there are a smooth vector field \(X\) and a constant \(\lambda\in\R\) such that
\[
\operatorname{Ric}(g)+\tfrac12\mathcal L_Xg=\lambda g.
\]
\textbf{Question 1.} For every initial metric on a general solvable Lie group, does a suitably normalized bracket flow converge to a bracket representing a Ricci soliton, and is the limit unique?

The supplied source does not specify the normalization, time reparametrization, convergence topology, or equivalence relation for uniqueness. In particular, it does not say whether uniqueness concerns a fixed initial metric or all initial metrics on the same group, or whether limits are identified under isometry, scaling, or changes of basis. It also does not require the limiting bracket to represent the original Lie group rather than a degeneration. Thus this question remains a formulation requiring those specifications; no particular normalization or uniqueness assertion is imposed here.

For the second question, a homogeneous Riemannian metric is one whose isometry group acts transitively. Its scalar curvature is spatially constant. A homogeneous Ricci-flow solution is a family of homogeneous metrics satisfying
\[
\frac{\partial g(t)}{\partial t}=-2\operatorname{Ric}(g(t)).
\]
\textbf{Question 2.} Suppose a homogeneous space \(M\) admits a homogeneous Riemannian metric \(h\) with \(\operatorname{scal}(h)>0\). Must every homogeneous Ricci-flow solution on \(M\) eventually have positive scalar curvature?

For a solution considered on its maximal forward interval \([0,T)\), where \(0<T\leq\infty\), a precise possible meaning of ``eventually'' is
\[
\exists\,t_0\in[0,T)\quad\forall\,t\in[t_0,T):
\operatorname{scal}(g(t))>0.
\]
This includes finite maximal existence times. The source does not explicitly prescribe this time convention. It also does not specify whether all metrics must be invariant under one fixed transitive group action on \(M=G/H\), or whether their transitive isometry groups may differ. These choices remain unresolved.
\subsection{Short English statement}
For arbitrary solvable Lie groups, ask whether suitably normalized bracket flow converges to a Ricci soliton with a unique limit. Separately, ask whether admitting a homogeneous metric of positive scalar curvature forces every homogeneous Ricci flow on the space eventually to have positive scalar curvature. The normalization, meaning of uniqueness, symmetry convention, and forward-time interpretation require clarification.
\subsection{Sources}
\begin{itemize}
\item Supplied local UnsolvedMath registry record, permanent ID \(30002691\), source problem code \texttt{catalogue entry}, title \emph{Convergence and Curvature of Homogeneous Ricci Flows}; collection: Oberwolfach Reports Open Problems. No individual proposer is supplied; the proposed year is recorded as 2014.
\item Supplied source citation: \emph{Mini-Workshop: Einstein Metrics, Ricci Solitons and Ricci Flow under Symmetry Assumptions} (2015), Oberwolfach Reports, pp.\ 2529--2568, DOI \texttt{10.4171/owr/2014/45}. \url{https://doi.org/10.4171/owr/2014/45}. The supplied literature list labels this source \emph{Oberwolfach Report 45/2014}; the citation's 2015 date and the report's 2014 designation are retained without reconciliation.
\item Supplied literature reference: \emph{On Homogeneous Ricci Solitons}. \url{https://doi.org/10.1093/qmath/hat028}.
\item Supplied literature reference: \emph{Scalar Curvature Behavior of Homogeneous Ricci Flows}. \url{https://arxiv.org/abs/1212.6558}.
\item Literature assessment dated exactly 2026-08-21: the supplied curation labels the problem partially solved and reports known bracket-flow fixed points and convergence results for important solvable and nilpotent classes, while describing general convergence, uniqueness, and the stated universal scalar-curvature behavior as unresolved. This is supplied dated curation, not an independently verified literature assessment; its status and cited claims require independent verification.
\end{itemize}
\end{document}
